% Generated from data/inputs.json raw_evidence; original and adopted values separated.
\begingroup\footnotesize\setlength{\tabcolsep}{3pt}\renewcommand{\arraystretch}{1.08}
\begin{longtable}{@{}>{\raggedright\arraybackslash}p{27mm}>{\raggedright\arraybackslash}p{62mm}>{\raggedright\arraybackslash}p{65mm}@{}}
\caption{参数证据与采用方法。页码均为本地PDF页序；原始值与本文预算分列。}\label{02_sram_dcim:tab:dcim-evidence}\\
\toprule 证据与定位 & 原值、单位和条件 & 采用值与换算理由\\\midrule\endfirsthead
\toprule 证据与定位 & 原值、单位和条件 & 采用值与换算理由\\\midrule\endhead
E01 \texttt{SDCIM-01}\newline p.1 Fig.1；p.2首段 & 物理阵列与完整INT8求值：128×128 bit；128×16 INT8；16项×16输出/拍；64 clock/VMM。28 nm，6T binary SRAM，8-bit signed/unsigned full-precision；输出23 bit & 一份128×128 INT8需8个容量分片；单分片64拍，顺序覆盖512拍。128个物理bit列除以每权重8 bit，仅16个逻辑输出；16项归约保留原结构\\\addlinespace[2pt]
E02 \texttt{SDCIM-01}\newline p.2 II-D；p.3 Fig.5 & 完整MAC槽与静态读：16 HCA +16 BFA；8次input-bit展开后换行组；23-bit RCA。静态dual-wordline access；每LBL 16 cell；BFA不在8输入位之间重切WL & 每轮含本地读保持、1-bit输入×8-bit权重、16项归约、符号/移位/累加；无独立read\_ns。完整计算周期覆盖读与算；64次行组选择支撑512个bit-round，不把静态WL重复解释成512次切换\\\addlinespace[2pt]
E03 \texttt{SDCIM-01}\newline p.4 Fig.10(a)；p.3 III & D6CIM计算时钟：0.9 V: 233 MHz；1.1 V: 正文360 MHz、图363 MHz。28 nm，VMM shmoo；温度未在该图明确报告；不采用最低0.6 V/30 MHz情景 & T\_C=4.3/5/10 ns；4.3 ns由1000/233向上取0.1 ns；5/10由共同数字节拍选择。0.9 V与共同近标称点相邻；不从1.1 V高速点外推0.9 V；三值为预算不是PVT分位数\\\addlinespace[2pt]
E04 \texttt{SDCIM-01}\newline p.1 Fig.1 & 普通SRAM写接口：Write Data[127:0]；Addr[6:0]；CEB/WEB。memory R/W column/row periphery，与VMM控制分开；未独立报告memory write ns & 128 bit=16个完整INT8权重/事务；单分片单行，单更新域。每权重8 binary cell；bitcell互补内部节点不产生额外逻辑payload；不把CIM频率当写频率\\\addlinespace[2pt]
E05 \texttt{CMOS-07}\newline p.2 Fig.2；p.9 Tables II/III & 普通SRAM完整同步周期锚点：455 MHz；48 bit×1k word；2RW 8T；1.05 V typical；25 °C读写功耗测试。28-nm eFlash工艺，高阈值；共同CLK上升沿捕获命令/地址/数据；图示连续write/write/read；频率非单独write-min测量 & 1000/455=2.197802… ns，取2.2 ns为写周期参考下端；5/10 ns为公共节拍预算。以完整时钟和同步写组织支持普通写，不用read access。跨48→128 bit与8T→6T是明确参考移植，需128路写驱动，非同芯片实测配对\\\addlinespace[2pt]
E06 \texttt{CMOS-07}\newline p.9 Table II & read access排除项：1.54 ns。28-nm 48-kbit 2RW SRAM，typical read access & 不用于t\_write或独立DCIM读附加项。read access只覆盖读访问，不是周期，更不是写完成时间\\\addlinespace[2pt]
E07 \texttt{CMOS-03}\newline p.1 II；p.3 IV-C & 6T写成功终点与周期交叉证据：50 FO4$\approx$1 ns；WL=25 FO4；15 fF/128 cell；SA offset0.1 V。HD 28 nm 6T，TT，nominal1 V；瞬态模型；周期末检查内部节点达到目标，支持back-to-back；目标BER<1e-9对应1MB/90\% yield & 仅支持同步完整写及数ns量级；不把1 ns当128-bit D6CIM实测写周期。约半周期WL脉冲$\ne$完整更新；需含节点翻转、保持/关WL和下一周期可读\\\addlinespace[2pt]
E08 \texttt{SDCIM-03}\newline p.1；p.2 Figs.11.7.2-5 & 其他28 nm数字组织：3 ns/clock@0.8 V；8 clock/8-bit乘法；$\Sigma$4/$\Sigma$9/$\Sigma$32 post-sum。dynamic logic；memory和CIM模式分离；lossless bitwise computation；不同post-sum组织 & 仅交叉支持数ns计算时钟和逐位完整精度；不把24 ns直接作为128×128 MVM。3 ns是本宏时钟；输出数、post-sum及本地接口与D6不同\\\addlinespace[2pt]
E09 共享基线\newline JSON共同条件/R0 & 共同逻辑/外围/边界：128×128 INT8；128×24-bit输出；T\_D=2/5/10 ns；R0=32项×16输出/轮。28 nm，0.9 V/25 °C参考，非固定PVT验证；单独立服务单元/更新域；无等面积约束 & 逻辑/精度/边界不变；D6型R2例外16项/轮；全向量捕获和提交2T\_D。R0=256轮是其资源选择；D6型512轮是16项而非32项的结果，不是共享计数错误\\\addlinespace[2pt]
E10 共享基线\newline 通用方法C及R4 & 完整写周期与控制覆盖：T\_front=[2+n\_b-$\chi$]T\_D；n\_b=$\chi$=1；完整周期已含则替代。普通同步memory端口，无另外定义的任务握手，边界在其命令/数据捕获到下一可计算周期 & 主情景Δ\_R=T\_M=2.2/5/10 ns；T\_front额外项为0；外层握手敏感性单列加2T\_D。完整memory周期已含命令解码、数据建立、写驱动和可复用终点；不会在同一阶段上再加T\_W\\\addlinespace[2pt]
\bottomrule\end{longtable}\endgroup
